\section{Research Gaps}

\subsection{Insufficient Practical Implementation of Secure P2P Rental Platforms}

While the body of academic literature on P2P vehicle sharing is substantial, it disproportionately favors theoretical models and prototype demonstrations over production-grade, deployable implementations. Systems such as eKar and Sharik, while cited as market examples, remain highly centralized in their trust and payment architectures, replicating the intermediary dependency they ostensibly seek to eliminate. There is a notable absence of rigorously documented full-stack implementations that address the practical engineering challenges of session security, payment localization, real-time communication, and dispute management in a coherent, integrated system. LuxeWay contributes to closing this gap by providing a fully deployed web platform with documented architectural decisions.

\subsection{Underexplored Hybrid CR+CS Models in the P2P Context}

The combined Car Rental and Car Sharing model has been explored primarily within B2C commercial fleet optimization contexts \cite{soppert2024benefit}. Its application to platforms where individual private owners act as supply-side participants introduces qualitatively different challenges: irregular vehicle availability, heterogeneous asset quality, variable owner responsiveness, and the absence of centralized fleet maintenance standards. These P2P-specific dynamics remain insufficiently modeled in the optimization literature, representing a significant gap between theoretical fleet economics and the operational realities of consumer-grade P2P marketplaces.

\subsection{Absence of UX-Centered System Design Documentation}

The engineering literature on vehicle rental platforms --- particularly within the blockchain and decentralized systems tradition --- demonstrates a persistent bias toward back-end security architecture at the expense of front-end usability and interaction design. Specific deficiencies include:

\begin{itemize}
    \item Incomplete documentation of UI component architecture and client-side state management strategies.
    \item Absence of accessibility compliance analysis for rental booking flows.
    \item Insufficient treatment of multi-language localization in markets such as Vietnam, where bilingual (Vietnamese/English) interfaces are commercially necessary.
    \item Minimal attention to responsive design across device form factors, despite the dominance of mobile browsing in Southeast Asian markets.
    \item Heavy dependence on theoretical models and simulations without large-scale real-world testing.
\end{itemize}

LuxeWay addresses these gaps by adopting a component-driven frontend architecture (React 18 + Radix UI + Tailwind CSS) with explicit internationalization support (i18next) and accessibility-first design patterns.

\subsection{Limitations of Subjective-Only Reputation Systems}

Voluntary review mechanisms remain the dominant trust-signaling instrument across existing platforms, despite well-documented vulnerabilities to fake review injection, rating reciprocity bias, and score convergence toward inflated means \cite{hu2019sharing}. The telematics-based approach proposed by Neifer et al.\ \cite{neifer2024trust} provides a compelling alternative but demands per-vehicle hardware investment that is impractical for casual private owners. The intermediate design space --- structured review decomposition, verified booking linkage, and temporal review analysis --- remains underexplored and represents an actionable research direction with immediate applicability to consumer P2P platforms.

\subsection{Unresolved Legal and Insurance Dimensions of P2P Rental}

The legal architecture supporting P2P vehicle sharing in Vietnam remains underdeveloped relative to mature markets such as the United States and the European Union. Key unresolved questions include the insurance liability boundary between personal vehicle policies and commercial rental usage, regulatory classification of P2P rental income, and the evidentiary weight of digitally generated rental contracts in dispute adjudication. Future iterations of LuxeWay must engage with these dimensions not merely as compliance requirements but as product design constraints that affect deposit structures, cancellation policies, and the scope of the platform's dispute resolution guarantee.
